% IEEE Journal Paper Template
% Federated continual learning: FedAvg vs dual-teacher (FLwF-2) vs clustered KD (FedKDC-CL).
% Flower (flwr) + PyTorch; MNIST/FMNIST/CIFAR-10 and MNIST grid sweeps.

\documentclass[journal]{IEEEtran}

\usepackage{cite}
\usepackage{amsmath,amssymb,amsfonts}
\usepackage{algorithmic}
\usepackage{algorithm}
\usepackage{graphicx}
\usepackage{textcomp}
\usepackage{xcolor}
\usepackage{booktabs}
\usepackage{multirow}
\usepackage{array}
\usepackage{url}
\usepackage{hyperref}
\usepackage{pifont}

\graphicspath{{images/}{./}}

\hypersetup{
    colorlinks=true,
    linkcolor=blue,
    citecolor=blue,
    urlcolor=blue
}

\def\BibTeX{{\rm B\kern-.05em{\sc i\kern-.025em b}\kern-.08em
    T\kern-.1667em\lower.7ex\hbox{E}\kern-.125emX}}

% All manuscript figures use a square ($1{:}1$) bounding box: \texttt{width} and
% \texttt{height} both set to \texttt{\string\linewidth} (relative units, not cm);
% the raster/vector aspect ratio is preserved inside the box (\texttt{keepaspectratio}).
\newcommand{\FigSquare}[1]{%
  \includegraphics[width=\linewidth,height=\linewidth,keepaspectratio]{#1}}
% Two square panels side-by-side in \texttt{figure*} (each panel $\approx 0.48\times$ column height).
\newcommand{\FigSqHalf}[1]{%
  \includegraphics[width=0.48\textwidth,height=0.48\textwidth,keepaspectratio]{#1}}

\begin{document}

\title{Can Federated Models Learn New Tasks Without Forgetting the Old Ones?}



\author{
Harshal Dharpure,
Anup Kumar Gond,
Riha Sanjay Kokode,
and Abdul Hadi%
\thanks{
Harshal Dharpure (2511AI30) is with the Department of Computer Science and Engineering,
Indian Institute of Technology Patna, Patna, Bihar, India
(e-mail: harshal\_2511ai30@iitp.ac.in).
}%
\thanks{
Anup Kumar Gond (2511CS01) is with the Department of Computer Science and Engineering,
Indian Institute of Technology Patna, Patna, Bihar, India
(e-mail: anup\_2511cs01@iitp.ac.in).
}%
\thanks{
Riha Sanjay Kokode (2511AI47) is with the Department of Computer Science and Engineering,
Indian Institute of Technology Patna, Patna, Bihar, India
(e-mail: riha\_2511ai47@iitp.ac.in).
}%
\thanks{
Abdul Hadi (2511MC17) is with the Department of Mathematics and Computing,
Indian Institute of Technology Patna, Patna, Bihar, India
(e-mail: hadi\_2511mc17@iitp.ac.in).
}
}

\maketitle

%------------------------------------------------------------------
\begin{abstract}
We present a reproducibility study of FLwF-2, a dual-teacher
knowledge-distillation approach to Federated Continual Learning, and
compare it to a FedAvg baseline. Both methods are implemented in
Flower~+~PyTorch and evaluated on MNIST, Fashion-MNIST and CIFAR-10
under a class-incremental two-task split, $K{\in}\{10,50,100\}$
clients, and five partition regimes (IID and Dirichlet
$\alpha{\in}\{0.01,0.1,0.5,1.0\}$). We track global accuracy/loss,
task-wise accuracy, forgetting, backward transfer and cumulative
communication cost. With the original FLwF-2 hyperparameters
($\alpha_{CE}{=}0.001$, $\beta_{KD}{=}0.7$, $T{=}2$), FedAvg almost
completely forgets Task 1 (forgetting $\approx 0.99$ on MNIST,
$0.80$ on CIFAR-10) while learning Task 2, whereas FLwF-2 preserves
Task 1 (forgetting $\approx 0$) but barely adapts to Task 2. The two
methods sit at opposite extremes of the stability--plasticity curve;
FLwF-2's forgetting profile is strictly better at zero extra
communication cost, validating the dual-teacher mechanism while
exposing a tunable knob that future work should adapt online.
Separately, we report a complementary Flower simulation study that
compares FedAvg to \emph{FedKDC-CL}---clustered federated knowledge
distillation with drift-aware scheduling---on MNIST under the same
$K$ and partition choices (IID and Dirichlet $\alpha$), using a fast
configuration aligned with our submission timeline; aggregate forgetting,
backward transfer and task-wise metrics appear in
Section~\ref{ssec:fedkdc_cl_results}.
\end{abstract}

\begin{IEEEkeywords}
Federated Learning, Continual Learning, Catastrophic Forgetting,
Knowledge Distillation, Dual-Teacher Distillation, FedAvg, FLwF-2,
FedKDC-CL, Clustered Distillation, Non-IID Data, Flower, PyTorch.
\end{IEEEkeywords}

%==================================================================
\section{Introduction}
\label{sec:intro}
%==================================================================

The proliferation of edge devices --- smartphones, wearables, IoT
sensors --- has created vast pools of locally generated data.
Centralizing this data for model training raises serious privacy,
legal and bandwidth concerns. Federated Learning (FL), introduced by
McMahan~et~al.~\cite{fedavg}, addresses this: clients train models
locally and share only parameters, which the server aggregates via
FedAvg without ever seeing raw data.

Two long-standing obstacles remain when FL is deployed in realistic
pervasive-computing settings:

\textbf{Problem 1 --- Catastrophic Forgetting.} Edge clients do not
receive a static stream of data. A wearable today might record
\textit{walking} activity, and tomorrow only \textit{sitting} data.
A neural network fine-tuned on the new task tends to dramatically
forget what it learned before~\cite{forgetting}. This effect compounds
in FL, because the server must simultaneously serve clients whose task
histories differ.

\textbf{Problem 2 --- Statistical Heterogeneity (non-IID).} Clients
hold non-identically distributed data: some clients see only digits
1 and 2, others only 7 and 8. Local gradients then point in different
directions, causing \textit{client drift} and degrading the global
model.

We focus on a particular instance of Federated Continual Learning
(FCL) that targets Problem 1 directly while still being subject to
Problem 2: the \emph{FLwF-2} method by Usmanova~et~al.~\cite{flwf},
in which each client uses two teachers --- its previous local model
and the previous server model --- to distill knowledge into the
current student during local training. We reproduce FLwF-2 in the
Flower framework, evaluate it against the FedAvg baseline across three
standard image datasets, and analyse its behaviour under varying
client counts and Dirichlet non-IID partitioning.

\textbf{Contributions.}
We provide an open Flower~\cite{flower}~+~PyTorch implementation of
FedAvg and FLwF-2 with a YAML-driven, resume-safe sweep harness, and
a 90-run empirical study (2 methods $\times$ 3 datasets $\times$ 3
client counts $\times$ 5 partitions) showing that --- with the
paper-recommended $(\alpha_{CE},\beta_{KD},T)=(0.001,0.7,2)$ ---
FLwF-2 essentially eliminates forgetting at the cost of plasticity,
while FedAvg exhibits the converse failure mode.
We additionally contribute an open implementation of FedKDC-CL in the
same codebase and a fifteen-configuration MNIST sweep (three client
counts $\times$ four Dirichlet alphas plus IID), summarised in
Tables~\ref{tab:fedkdc_grid_dirichlet}--\ref{tab:fedkdc_extra_runs},
which isolates the clustered-KD baseline against FedAvg under identical
partitioning without mixing task streams across datasets.

The remainder of the paper is organised as follows.
Section~\ref{sec:background} reviews FL, knowledge distillation and
non-IID partitioning. Section~\ref{sec:method} formalises the
class-incremental FCL setting and defines the FLwF-2 dual-teacher
loss exactly as implemented. Section~\ref{sec:setup} describes the
experimental harness. Section~\ref{sec:results} presents the
quantitative results (including FedKDC-CL vs.\ FedAvg in
Section~\ref{ssec:fedkdc_cl_results}) and Section~\ref{sec:analysis}
the in-depth analysis. Section~\ref{sec:limitations} discusses limitations and
future directions and Section~\ref{sec:conclusion} concludes.
Appendix~\ref{app:results} tabulates the full FLwF-2 sweep.
FedKDC-CL grid hyperparameters are summarised in Section~\ref{ssec:fedkdc_cl_results}.

%==================================================================
\section{Background}
\label{sec:background}
%==================================================================

\subsection{Federated Averaging (FedAvg)}

A server coordinates $K$ clients with private datasets $D_k$. At
each communication round $r$, a fraction $\rho$ of clients are
sampled; each selected client downloads the current global model
$\theta^{r}$, runs $E$ local epochs of SGD on $D_k$, and returns
$\theta_k^{r+1}$. The server aggregates:
\begin{equation}
    \theta^{r+1} = \sum_{k\in S_r} \frac{|D_k|}{\sum_{j\in S_r}|D_j|}\,
    \theta_k^{r+1}.
\end{equation}
We use this as our baseline throughout.

\subsection{Knowledge Distillation}

Knowledge Distillation (KD)~\cite{hinton_kd} transfers knowledge from a
teacher to a student via soft predictions. With temperature $T$:
\begin{equation}
    \pi_i(x) = \frac{\exp(o_i(x)/T)}{\sum_j \exp(o_j(x)/T)},
\end{equation}
and the KD objective is
\begin{equation}
    \mathcal{L}_{\textrm{KD}} = T^2\,
    \mathrm{KL}\!\left(\sigma\!\left(\tfrac{o^{S}(x)}{T}\right)\;\big\|\;
    \sigma\!\left(\tfrac{o^{T}(x)}{T}\right)\right),
\end{equation}
where $o^{S}, o^{T}$ are student and teacher logits and $\sigma$ is
softmax. We follow the standard PyTorch idiom of computing
\texttt{log\_softmax} on the student and \texttt{softmax} on the
teacher and feeding both into \texttt{KLDivLoss(reduction="batchmean")}.

\subsection{Non-IID Partitioning via Dirichlet}

To control statistical heterogeneity, we sample per-class client
weights from a Dirichlet distribution $\mathrm{Dir}(\alpha)$ with a
single concentration parameter $\alpha$. As $\alpha\!\to\!0$, each
client receives a near-singleton class set (extreme non-IID); as
$\alpha\!\to\!\infty$, the partition tends to IID. We evaluate
$\alpha\in\{0.01, 0.1, 0.5, 1.0\}$ plus a uniformly random IID split.

\subsection{Federated Continual Learning Setting}

We adopt a class-incremental FCL setting with $T{=}2$ tasks. For each
dataset, classes are split into a first half (Task~1) and a second
half (Task~2) --- e.g.\ for CIFAR-10, Task~1 is classes $\{0,1,2,3,4\}$
and Task~2 is classes $\{5,6,7,8,9\}$. The federation first runs for
$R_1$ communication rounds on data restricted to Task~1, then for
$R_2$ rounds on data restricted to Task~2. Tasks are never mixed.
This is the simplest non-trivial FCL scenario in which catastrophic
forgetting is easy to observe and measure.

\subsection{Related Federated and Continual-Learning Context}

Class-incremental learning under covariate or task shift has been studied extensively in the centralized setting via rehearsal buffers, regularisation toward previous weights, and distillation from frozen snapshots (cf.\ learning without forgetting~\cite{lwf}). Federated optimisation approaches such as FedProx~\cite{fedprox} address statistical heterogeneity by proximal regularisation to the global iterate rather than explicit continual-learning objectives; they do not, by themselves, resolve forgetting across task boundaries. Ensemble distillation at the server (FedDF~\cite{feddf}) and peer knowledge matching (FedMD~\cite{fedmd}) reduce communication of raw logits but target cross-client fusion rather than sequential tasks on each device. Our FLwF-2 reproduction anchors one design point (dual teachers per client), while the FedKDC-CL extension explores clustered teachers under the same partition taxonomy---both complement heterogeneity-aware optimisation lines without replacing them.

%==================================================================
\section{Method}
\label{sec:method}
%==================================================================

\subsection{FLwF-2: Dual-Teacher Distillation}
\label{ssec:flwf2}

Each client maintains, in addition to the current model $\theta_k$
trained as the \emph{student}, two teachers:

\begin{itemize}
    \item \textbf{Teacher 1 (\emph{prev\_model})}: a frozen snapshot of
        the client's local model at the end of the \emph{previous} task.
        It is initialised from the server model on the first round of
        the experiment and is updated only at task boundaries (i.e.\ at
        the end of Task~1).
    \item \textbf{Teacher 2 (\emph{server\_model})}: a frozen snapshot
        of the current global model received from the server at the
        start of the current round. It carries knowledge from the
        entire federation.
\end{itemize}

\noindent
The local objective combines a cross-entropy term and two distillation
terms:
\begin{equation}
\boxed{\;
\mathcal{L}_{\mathrm{FLwF\text{-}2}}
   = \alpha_{CE}\,\mathcal{L}_{CE}
   + \beta_{KD}\,\mathcal{L}_{KD}^{\text{prev}}
   + (1\!-\!\alpha_{CE}\!-\!\beta_{KD})\,
     \mathcal{L}_{KD}^{\text{server}},
\;}
\label{eq:flwf2-loss}
\end{equation}
where $\mathcal{L}_{KD}^{\text{prev}}$ and
$\mathcal{L}_{KD}^{\text{server}}$ are KL-divergence terms between the
temperature-scaled student logits and each of the two frozen teachers.
Following the original paper, we use $\alpha_{CE}{=}0.001$,
$\beta_{KD}{=}0.7$ and $T{=}2$.

\subsection{Algorithm}

Algorithm~\ref{alg:flwf2} summarises one local training step under
FLwF-2 as implemented in our Flower client.

\begin{algorithm}[t]
\caption{FLwF-2 local training (one batch).}
\label{alg:flwf2}
\begin{algorithmic}[1]
\REQUIRE batch $(x,y)$, student $\theta$, prev teacher $\theta^{\text{prev}}$, server teacher $\theta^{\text{srv}}$, hyperparameters $\alpha_{CE},\beta_{KD},T$.
\STATE $z_S \leftarrow f_\theta(x)$ \COMMENT{student logits}
\STATE $z_P \leftarrow f_{\theta^{\text{prev}}}(x).\mathrm{detach}()$
\STATE $z_V \leftarrow f_{\theta^{\text{srv}}}(x).\mathrm{detach}()$
\STATE $\mathcal{L}_{CE} \leftarrow \mathrm{CE}(z_S, y)$
\STATE $\mathcal{L}_{KD}^{\text{prev}} \leftarrow T^2 \cdot \mathrm{KL}\!\left(\log\sigma(z_S/T)\;\|\;\sigma(z_P/T)\right)$
\STATE $\mathcal{L}_{KD}^{\text{srv}}  \leftarrow T^2 \cdot \mathrm{KL}\!\left(\log\sigma(z_S/T)\;\|\;\sigma(z_V/T)\right)$
\STATE $\mathcal{L} \leftarrow \alpha_{CE}\mathcal{L}_{CE} + \beta_{KD}\mathcal{L}_{KD}^{\text{prev}} + (1{-}\alpha_{CE}{-}\beta_{KD})\mathcal{L}_{KD}^{\text{srv}}$
\STATE backprop $\mathcal{L}$ through $\theta$ only
\STATE \textbf{at task end:} $\theta^{\text{prev}} \leftarrow \theta$
\end{algorithmic}
\end{algorithm}

The \texttt{prev\_model} update at task boundaries (Algorithm~\ref{alg:flwf2},
last line) is signalled by an explicit \texttt{task\_end} flag passed
from server to client via Flower's \texttt{config} dictionary.

\subsection{Metrics}
\label{ssec:metrics}

We log the following each communication round:

\begin{itemize}
    \item \textbf{Global Test Accuracy / Loss} on the full 10-class
        test set (same set across all clients and rounds).
    \item \textbf{Task-wise Accuracy} $a_t$ for each task $t\!\in\!\{1,2\}$,
        i.e.\ the accuracy restricted to that task's classes.
    \item \textbf{Catastrophic Forgetting} for task $t$:
        $F_t = \max_{r\le R_t}\,a_t^{(r)} - a_t^{(R)}$,
        where $R$ is the final round and $R_t$ is the last round in
        which the client trained on task $t$.
    \item \textbf{Backward Transfer (BWT)}: the change in accuracy on
        Task~1 from the end of Task~1 to the end of Task~2, signed.
    \item \textbf{Convergence Round}: the smallest round at which
        global test accuracy first reaches $80\%$.
    \item \textbf{Cumulative Communication Cost (bytes)}:
        $\sum_r (\textit{num\_sampled\_clients}_r) \cdot 2 \cdot
        (\textit{model size in bytes})$ --- both directions, full
        weights, FP32 (no compression).
\end{itemize}

%==================================================================
\section{Experimental Setup}
\label{sec:setup}
%==================================================================

\subsection{Framework, Reproducibility and Hardware}

All experiments use Flower (\texttt{flwr})~\cite{flower} with the
\texttt{simulation} extra (Ray-backed virtual clients) and PyTorch
$\geq 2.0$ as the deep learning backend. Reproducibility is enforced
by fixing \texttt{random.seed(42)},
\texttt{numpy.random.seed(42)} and \texttt{torch.manual\_seed(42)} in
both the server and every client. Each (method, dataset, clients,
partition) combination is executed in a fresh subprocess so that Ray
state, CUDA caches and dataloader workers are released cleanly between
runs. Sweeps run on a shared server with NVIDIA A100 GPUs.

\subsection{Model}

Both methods use the same lightweight CNN
\textit{Conv(32) $\to$ ReLU $\to$ MaxPool $\to$ Conv(64) $\to$ ReLU
$\to$ MaxPool $\to$ FC $\to$ Softmax}, with input bilinearly
interpolated to $32{\times}32$ (so that MNIST/FMNIST and CIFAR-10
share a single architecture); the only difference is the input channel
count (1 vs.\ 3). The total trainable model size is $\approx
60\,654$~parameters $\to 241.4$~KB at FP32, which determines the
per-round communication payload.

\subsection{Datasets and Partitioning}

\begin{itemize}
    \item \textbf{MNIST}: 60{,}000 training / 10{,}000 test images,
        10 classes, $28{\times}28$ grayscale.
    \item \textbf{Fashion-MNIST (FMNIST)}: 60{,}000 / 10{,}000,
        10 classes, $28{\times}28$ grayscale.
    \item \textbf{CIFAR-10}: 50{,}000 / 10{,}000, 10 classes,
        $32{\times}32$ RGB.
\end{itemize}

For each dataset we form a class-incremental two-task split
(classes $\{0,1,2,3,4\}$ for Task~1 and $\{5,6,7,8,9\}$ for Task~2).
Per-task data is then partitioned across clients in two ways:
(i)~uniform IID and (ii)~$\mathrm{Dir}(\alpha)$ with
$\alpha\in\{0.01,0.1,0.5,1.0\}$.
A re-balancing pass guarantees a minimum of 1 sample per (client,
task) so that no client gets an empty Task~1 or Task~2 dataset, even
under extreme heterogeneity.

\subsection{Federated and Training Configuration}

% Universal training and federated configuration.
\begin{table}[h]
\centering
\caption{Universal training and federated configuration.}
\label{tab:config}
\renewcommand{\arraystretch}{1.15}
\begin{tabular}{ll}
\toprule
\textbf{Parameter} & \textbf{Value} \\
\midrule
Algorithm                          & FedAvg / FLwF-2 \\
Clients $K$                        & 10, 50, 100 \\
Client fraction per round $\rho$   & 0.5 \\
Rounds Task 1 / Task 2             & 10 / 10 (total 20) \\
Local epochs                       & 5 \\
Local batch size                   & 32 \\
Optimizer                          & SGD, momentum $0.9$ \\
Learning rate                      & 0.01 \\
FLwF-2 $(\alpha_{CE},\beta_{KD},T)$ & $(0.001,\,0.7,\,2)$ \\
\bottomrule
\end{tabular}
\end{table}


\subsection{Sweep}

The full grid is $2~\textrm{methods} \times 3~\textrm{datasets}
\times 3~\textrm{client counts} \times 5~\textrm{partitions} =
90~\textrm{configurations}$. Each configuration writes per-run logs that are merged into a single
aggregate metrics table. The sweep is resume-safe and supervised to
tolerate transient Ray failures.

\subsection{Logging, Artifacts, and Third-Party Data}

Every subprocess records structured metrics for the FLwF-2 sweep and for
the FedKDC-CL launcher grid, plus high-resolution plots where generated.
Public benchmarks (MNIST, FMNIST, CIFAR-10) are accessed through
Flower-compatible dataset loaders under their respective licences; no
personally identifiable information is processed. Seeds are fixed so
tables and figures are reproducible modulo GPU nondeterminism for CUDA
kernels.

%==================================================================
\section{Experimental Results}
\label{sec:results}
%==================================================================

\subsection{Headline Comparison}

Table~\ref{tab:headline} contrasts FedAvg and FLwF-2 on representative
settings ($K{=}10$ clients, IID and $\mathrm{Dir}(0.1)$, all three
datasets). Per-task accuracy is reported as the global test accuracy
restricted to that task's classes. ``T1$\to$T1'' is Task~1 accuracy
\emph{at the end of Task~1}; ``T1$\to$T2'' is Task~1 accuracy
\emph{at the end of Task~2} (i.e., after the model has been trained on
Task~2 for $R_2$ rounds); ``T2$\to$T2'' is Task~2 accuracy at the end
of Task~2. \emph{Forgetting} is the drop in Task~1 accuracy from the
end of Task~1 to the end of Task~2.

% Headline comparison (final round; K=10 clients).
\begin{table*}[t]
\centering
\caption{Headline comparison (final round; $K{=}10$ clients).
``Avg'' is the unweighted mean of T1$\to$T2 and T2$\to$T2.
Forgetting is for Task~1.}
\label{tab:headline}
\renewcommand{\arraystretch}{1.15}
\begin{tabular}{llcccccc}
\toprule
\textbf{Method} & \textbf{Dataset} & \textbf{Partition}
 & \textbf{T1$\to$T1 (\%)} & \textbf{T1$\to$T2 (\%)}
 & \textbf{T2$\to$T2 (\%)} & \textbf{Avg (\%)}
 & \textbf{Forgetting} \\
\midrule
FedAvg & MNIST    & IID       &  99.7 &   0.0 &  99.4 &  49.7 & 1.00 \\
FedAvg & MNIST    & Dir(0.1)  &  98.4 &   0.0 &  98.0 &  49.0 & 0.99 \\
FLwF-2 & MNIST    & IID       &  99.7 &  99.8 &   0.0 &  49.9 & \textbf{0.00} \\
FLwF-2 & MNIST    & Dir(0.1)  &  98.3 &  98.3 &   0.0 &  49.2 & \textbf{0.01} \\
\midrule
FedAvg & FMNIST   & IID       &  93.5 &   0.0 &  97.7 &  48.9 & 0.94 \\
FedAvg & FMNIST   & Dir(0.1)  &  73.0 &   0.0 &  95.2 &  47.6 & 0.83 \\
FLwF-2 & FMNIST   & IID       &  93.5 &  93.4 &   0.0 &  46.7 & \textbf{0.00} \\
FLwF-2 & FMNIST   & Dir(0.1)  &  73.8 &  73.9 &   0.0 &  37.0 & \textbf{0.09} \\
\midrule
FedAvg & CIFAR-10 & IID       &  79.9 &   0.0 &  86.5 &  43.2 & 0.80 \\
FedAvg & CIFAR-10 & Dir(0.1)  &  50.3 &   0.0 &  72.6 &  36.3 & 0.60 \\
FLwF-2 & CIFAR-10 & IID       &  80.2 &  80.2 &   0.0 &  40.1 & \textbf{0.00} \\
FLwF-2 & CIFAR-10 & Dir(0.1)  &  54.3 &  54.4 &   0.0 &  27.2 & \textbf{0.06} \\
\bottomrule
\end{tabular}
\end{table*}


The two methods exhibit inverted failure modes: under FedAvg
T1$\to$T2 collapses to $0\%$ in every configuration while Task~2 is
learned cleanly; under FLwF-2 Task~1 is fully preserved
(T1$\to$T2 $\approx$ T1$\to$T1, forgetting $\approx 0$) but Task~2
stays at $\approx 0\%$. We analyse this tradeoff in detail in
Section~\ref{sec:analysis}.

\subsection{Effect of Non-IID Heterogeneity}

Figure~\ref{fig:iid_vs_noniid} reports final-round global test
accuracy across all five partition regimes for both methods,
aggregated over all (dataset, client-count) combinations. The IID
setting forms the upper envelope; pushing $\alpha$ towards
$0.01$ steadily erodes accuracy. In this aggregated view, FLwF-2 is
slightly more robust than FedAvg in the moderate non-IID regime
($\alpha\in\{0.1,0.5,1.0\}$), reflecting the regularising effect of
the server-side teacher.

\begin{figure}[t]
\centering
\FigSquare{iid_vs_noniid_final_accuracy.png}
\caption{Final-round global test accuracy for FedAvg and FLwF-2 under
all five partition regimes (IID and $\mathrm{Dir}(\alpha)$,
$\alpha\in\{0.01,0.1,0.5,1.0\}$), aggregated across all
(dataset, $K$) combinations. Bars are means; error bars are dropped
for clarity.}
\label{fig:iid_vs_noniid}
\end{figure}

\subsection{Convergence and Forgetting over Rounds}

Figures~\ref{fig:acc_rounds}--\ref{fig:forgetting_rounds} plot global
accuracy, loss and the Task~1 forgetting metric versus round. All
three show the same picture: both methods climb to a Task-1 plateau
($\approx 49\%$ on MNIST/FMNIST, the upper bound when only Task~1
classes are testable) in the first ten rounds; at round 11 the task
switches and FedAvg's forgetting jumps to $\sim\!1$ within a few
rounds, whereas FLwF-2 keeps it pinned near zero for the rest of the
run. Figure~\ref{fig:method_box} confirms this as a per-method
distribution: FLwF-2's stability prior bounds its variance, while
FedAvg occasionally collapses below $20\%$ in the most extreme
non-IID configurations.

\begin{figure}[t]
\centering
\FigSquare{accuracy_vs_rounds.png}
\caption{Global test accuracy versus communication rounds across all
runs (lines coloured by method, styled by partition). The drop at
round 11 is the Task~1 $\to$ Task~2 boundary.}
\label{fig:acc_rounds}
\end{figure}

\begin{figure}[t]
\centering
\FigSquare{loss_vs_rounds.png}
\caption{Global test loss versus rounds across all runs.}
\label{fig:loss_rounds}
\end{figure}

\begin{figure}[t]
\centering
\FigSquare{forgetting_vs_rounds.png}
\caption{Catastrophic-forgetting metric for Task~1 versus round.
FedAvg climbs to $\approx 1$ shortly after the task switch; FLwF-2
keeps it near $0$.}
\label{fig:forgetting_rounds}
\end{figure}

\begin{figure}[t]
\centering
\FigSquare{fedavg_vs_flwf2_final_accuracy.png}
\caption{Distribution of final-round global accuracy across the full
sweep, per method.}
\label{fig:method_box}
\end{figure}

\subsection{Communication Cost}

Both methods exchange the same FP32 weights, so cumulative
communication cost depends only on model size, $K$, and number of
rounds. With our 60{,}654-parameter CNN ($\approx 241.4$~KB/message),
20 rounds and $\rho{=}0.5$, totals are $47.8$, $239.1$ and
$478.3$~MB for $K{=}10,50,100$ respectively (full per-run values in
Table~\ref{tab:full_results}). FLwF-2 derives both teachers locally
and therefore adds no communication overhead.

\subsection{FedKDC-CL vs.\ FedAvg (MNIST Simulation Study)}
\label{ssec:fedkdc_cl_results}

In parallel to the FLwF-2 study above, we implemented and swept a
second federated continual-learning baseline comparison in the same
Flower~+~PyTorch stack: \textbf{FedAvg} versus \textbf{FedKDC-CL}, a
clustered knowledge-distillation variant with drift-aware scheduling
and entropy-aware KD temperature. Experiments use the same MNIST class split
(Task~1: $\{0,\ldots,4\}$, Task~2: $\{5,\ldots,9\}$), Dirichlet and IID
partitions, and $K\!\in\!\{10,50,100\}$ clients with fraction $\rho{=}0.5$.
The FedKDC-CL grid uses a shortened horizon for throughput: two communication
rounds per task; SGD with learning rate $0.01$, batch size $32$, momentum
$0.9$, two local epochs per round; random seed~$42$; two clusters; KD enabled
with one KD epoch per KD phase, equal CE/KD weights, temperature in $[1,6]$;
drift guard with promotion every two rounds; Ray virtual clients on CPU with
four CPUs per client and centralized evaluation where accelerators are
available. Metrics are summary aggregates over the full ten-class label space: \emph{final average accuracy},
\emph{average forgetting}, and \emph{backward transfer} (BWT), comparing both
methods for each configuration.

Table~\ref{tab:fedkdc_grid_dirichlet} lists all four Dirichlet
concentrations $\alpha\!\in\!\{0.01,0.1,0.5,1.0\}$; Table~\ref{tab:fedkdc_grid_iid}
lists IID runs for each $K$.
The full launcher completed fifteen configurations successfully on 2026-05-09.
Table~\ref{tab:fedkdc_extra_runs} summarises additional pilot runs from
earlier development.

% FedKDC-CL fast grid: Dirichlet non-IID.
\begin{table*}[t]
\centering
\footnotesize
\caption{FedKDC-CL vs.\ FedAvg on MNIST continual learning (two class-incremental tasks; summary metrics over all classes).
Non-IID partitions via $\mathrm{Dir}(\alpha)$.
Setup: two communication rounds per task, client fraction $0.5$, virtual clients on CPU, seed $42$
(details in Section~\ref{ssec:fedkdc_cl_results}).
\emph{FAvgAcc}/\emph{FFrg}/\emph{FBWT}: FedAvg final average accuracy, average forgetting, backward transfer;
\emph{KKDC}/\emph{KFrg}/\emph{KBWT}: FedKDC-CL.}
\label{tab:fedkdc_grid_dirichlet}
\renewcommand{\arraystretch}{1.12}
\begin{tabular}{@{}c l cccccc@{}}
\toprule
$K$ & $\alpha$ & \textbf{FAvgAcc} & \textbf{FFrg} & \textbf{FBWT}
& \textbf{KKDC} & \textbf{KFrg} & \textbf{KBWT} \\
\midrule
10 & 0.01 & 0.1706 & 0.1993 & -0.3985 & 0.1897 & 0.0955 & -0.1911 \\
10 & 0.1 & 0.2408 & 0.4774 & -0.9549 & 0.3414 & 0.4866 & -0.9731 \\
10 & 0.5 & 0.4200 & 0.4805 & -0.9611 & 0.4278 & 0.4869 & -0.9737 \\
10 & 1.0 & 0.4742 & 0.4927 & -0.9854 & 0.4638 & 0.4910 & -0.9819 \\
50 & 0.01 & 0.0964 & 0.0392 & -0.0784 & 0.1004 & 0.0000 & 0.1868 \\
50 & 0.1 & 0.2330 & 0.2435 & -0.4871 & 0.2363 & 0.2972 & -0.5945 \\
50 & 0.5 & 0.4416 & 0.4606 & -0.9212 & 0.4200 & 0.4648 & -0.9296 \\
50 & 1.0 & 0.4356 & 0.4784 & -0.9568 & 0.4284 & 0.4560 & -0.9120 \\
100 & 0.01 & 0.0964 & 0.0392 & -0.0784 & 0.1004 & 0.0000 & 0.1868 \\
100 & 0.1 & 0.2696 & 0.1242 & -0.2485 & 0.1324 & 0.2729 & -0.5458 \\
100 & 0.5 & 0.2637 & 0.4419 & -0.8838 & 0.3110 & 0.3793 & -0.7585 \\
100 & 1.0 & 0.3261 & 0.4313 & -0.8626 & 0.3034 & 0.3709 & -0.7418 \\
\bottomrule
\end{tabular}
\end{table*}

\begin{table}[t]
\centering
\footnotesize
\caption{FedKDC-CL vs.\ FedAvg under IID partitioning (same setup as Table~\ref{tab:fedkdc_grid_dirichlet}).}
\label{tab:fedkdc_grid_iid}
\renewcommand{\arraystretch}{1.12}
\begin{tabular}{@{}c cccccc@{}}
\toprule
$K$ & \textbf{FAvgAcc} & \textbf{FFrg} & \textbf{FBWT} & \textbf{KKDC} & \textbf{KFrg} & \textbf{KBWT} \\
\midrule
10 & 0.4853 & 0.4916 & -0.9833 & 0.4770 & 0.4891 & -0.9782 \\
50 & 0.4342 & 0.4781 & -0.9562 & 0.3912 & 0.4674 & -0.9348 \\
100 & 0.3548 & 0.4466 & -0.8932 & 0.3038 & 0.3871 & -0.7743 \\
\bottomrule
\end{tabular}
\end{table}
% Earlier FedKDC-CL-related pilots (development).
\begin{table}[t]
\centering
\footnotesize
\caption{Additional FedKDC-CL pilots (development and smoke tests; MNIST). Same aggregate definitions as Tables~\ref{tab:fedkdc_grid_dirichlet}--\ref{tab:fedkdc_grid_iid}.}
\label{tab:fedkdc_extra_runs}
\renewcommand{\arraystretch}{1.12}
\begin{tabular}{@{}llccc@{}}
\toprule
\textbf{Pilot} & \textbf{Method} & \textbf{Acc.} & \textbf{Forg.} & \textbf{BWT} \\
\midrule
\multirow{2}{*}{P1}
 & FedAvg & 0.0964 & 0.0392 & -0.0784 \\
 & FedKDC-CL & 0.1004 & 0.0000 & 0.1868 \\
\midrule
\multirow{2}{*}{P2}
 & FedAvg & 0.1988 & 0.4896 & -0.9792 \\
 & FedKDC-CL & 0.1800 & 0.4929 & -0.9858 \\
\midrule
P3 & FedKDC-CL & 0.1654 & 0.1004 & -0.2008 \\
\bottomrule
\end{tabular}
\end{table}


Figure~\ref{fig:fedkdc_compare} shows a representative task-wise accuracy
comparison for IID with $K{=}100$;
Figure~\ref{fig:fedkdc_rounds_grid} shows per-task accuracy versus round for
the same setting (FedAvg vs.\ FedKDC-CL).
Qualitatively, FedKDC-CL and
FedAvg track closely under this fast configuration because both methods share
the same local optimiser and only diverge on alternating KD rounds for
FedKDC-CL. Interpreting Tables~\ref{tab:fedkdc_grid_dirichlet}--\ref{tab:fedkdc_grid_iid}
together, larger $\alpha$ and IID partitions generally raise final average
accuracy (less extreme label skew), while very small $\alpha$ paired with large
$K$ produces sparse class coverage per client and drives both methods toward
degenerate summary aggregates. Backward transfer remains strongly negative in
most cells, consistent with catastrophic forgetting on Task~1 after Task~2
training under short horizons.

\begin{figure}[t]
\centering
\FigSquare{fedkdc/comparison_acc_vs_tasks.png}
\caption{Representative FedKDC-CL vs.\ FedAvg comparison: final accuracy per
continual task on MNIST (IID, $K{=}100$).}
\label{fig:fedkdc_compare}
\end{figure}

\begin{figure*}[t]
\centering
\FigSqHalf{fedkdc/fedavg_mnist_0_4_acc_vs_rounds.png}\hfill
\FigSqHalf{fedkdc/fedkdc_mnist_0_4_acc_vs_rounds.png}\\[0.35em]
{\footnotesize Task~1 (classes 0--4):\quad FedAvg (left),\quad FedKDC-CL (right)}\\[0.75em]
\FigSqHalf{fedkdc/fedavg_mnist_5_9_acc_vs_rounds.png}\hfill
\FigSqHalf{fedkdc/fedkdc_mnist_5_9_acc_vs_rounds.png}\\[0.35em]
{\footnotesize Task~2 (classes 5--9):\quad FedAvg (left),\quad FedKDC-CL (right)}
\caption{Per-task accuracy vs.\ communication round (IID, $K{=}100$).
Row~1: Task~1; row~2: Task~2. Left: FedAvg; right: FedKDC-CL.}
\label{fig:fedkdc_rounds_grid}
\end{figure*}

%==================================================================
\section{Analysis}
\label{sec:analysis}
%==================================================================

\subsection{The Stability--Plasticity Tradeoff}

Equation~\ref{eq:flwf2-loss} weights cross-entropy by $\alpha_{CE}{=}0.001$
and the prior-knowledge KD term by $\beta_{KD}{=}0.7$. The new-task
gradient signal carries effective weight only $0.001$, while the two
teachers --- both unaware of Task~2 --- jointly carry weight $0.999$.
With this configuration, FLwF-2 has essentially no incentive to deviate
from the previous model on Task-2 batches, which is exactly what we
observe: Task~2 accuracy stays near $0\%$ throughout rounds 11--20.
This is not an implementation bug; it is the predicted equilibrium of
the loss with these hyperparameters and a class-incremental split.
The original paper anticipated this and proposed a
\emph{FLwF-2 + Fine-Tuning} hybrid that switches the loss weighting
adaptively. We deliberately stayed with the literal hyperparameters to
isolate the dual-teacher contribution.

\subsection{Effect of Client Count}

Across all three datasets, increasing $K$ from 10 to 100 (with
$\rho{=}0.5$ fixed, so the number of sampled clients per round grows
linearly) produces a small accuracy drop, primarily because each
client's local dataset shrinks and per-client gradient noise
increases. The effect is mild on MNIST and FMNIST (well within
$\pm 2\%$ for both methods at IID) and more pronounced on CIFAR-10
($\approx 5{-}6\%$ between $K{=}10$ and $K{=}100$).

\subsection{Effect of Dirichlet $\alpha$}

For $\alpha\!\in\!\{0.5,1.0\}$ and IID, both methods reach near-clean
Task accuracy. At $\alpha\!=\!0.1$ both still converge but to a lower
plateau. At $\alpha\!=\!0.01$, label distribution is essentially
single-class per client; both methods collapse to $\approx 10\%$ on
MNIST/CIFAR-10 and roughly fail to separate the classes. FLwF-2
inherits this failure mode --- the dual-teacher prior cannot rescue an
underlying federation that has no consistent gradient signal.

\subsection{Backward Transfer and Convergence Round}

Because FLwF-2 freezes most of the prior representation, its backward
transfer on Task~1 is essentially zero ($\le 0.5\%$ change after
Task~2), whereas FedAvg's BWT is strongly negative ($\approx -1.0$ on
MNIST), as expected.

The course-mandated convergence-round metric (smallest round at which
\emph{global} accuracy first reaches $80\%$) is undefined in our
class-incremental setting: only one half of the label space is
testable at a time, so the global ceiling is $\approx 50\%$ by
construction and no run attains $80\%$. We therefore report ``n/a''
in Table~\ref{tab:full_results} rather than fabricate a value;
redefining the metric per active task is left to future work.

\subsection{Cross-Reading FLwF-2 and FedKDC-CL Metrics}

The FLwF-2 experiments prioritise per-task accuracy paths and forgetting on
three datasets with twenty communication rounds per task, whereas the FedKDC-CL
grid intentionally shortens training (two rounds per task) to obtain a dense
coverage of $(K,\alpha)$ settings. Readers should therefore compare magnitudes
only within each study, not across tables: absolute summary aggregates are not
calibrated between the two protocols. Conceptually, FLwF-2 foregrounds
dual-server-and-local teachers with extreme hyperparameter-induced plasticity
collapse, while FedKDC-CL foregrounds periodic KD phases driven by cluster
centroids and drift promotion---both exhibit stability--plasticity tension but
through incomparable loss schedules.

%==================================================================
\section{Limitations and Future Work}
\label{sec:limitations}
%==================================================================

\paragraph{Limitations.}
We use only two tasks, the literal $(\alpha_{CE},\beta_{KD})$ from
the original paper, and a homogeneous CNN. We do not compare against
FLwF-1 or the FLwF-2 + Fine-Tuning hybrid.
The FedKDC-CL sweep uses MNIST only (fast grid), omits FMNIST/CIFAR-10
stress tests, and relies on CPU-side virtual clients for numerical
stability---these choices bias toward sweep throughput rather than minimising
single-run latency on accelerators.
Simulation-scale findings may differ from cross-device deployments where
stragglers, partial participation, and asynchronous uploads violate our
synchronous round assumptions.

\paragraph{Future Work.}
The most direct next step is an adaptive
$(\alpha_{CE},\beta_{KD})$ schedule that interpolates between the
plasticity and stability extremes (e.g.\ start plastic, anneal
towards stability), or equivalently the FLwF-2 + Fine-Tuning hybrid
proposed in~\cite{flwf}. Beyond that: longer task sequences on
CIFAR-100, soft-label-only exchange à la FedMD~\cite{fedmd}~/
FedDF~\cite{feddf} for architecture-heterogeneous federations, and
a privacy analysis quantifying any incremental leakage from sharing
weights versus soft labels.

%==================================================================
\section{Conclusion}
\label{sec:conclusion}
%==================================================================

A 90-run sweep of FedAvg and FLwF-2 in Flower~+~PyTorch confirms that
dual-teacher distillation eliminates catastrophic forgetting at zero
extra communication cost, but with the literal
$(\alpha_{CE},\beta_{KD})=(0.001,0.7)$ it does so at the cost of
near-zero plasticity --- the mirror image of FedAvg's high-plasticity,
high-forgetting failure mode. The natural next step is an adaptive
scheme that interpolates between these extremes; our resume-safe
sweep harness is designed to make such follow-up experiments
straightforward.
The supplementary FedKDC-CL MNIST grid (Section~\ref{ssec:fedkdc_cl_results})
provides an orthogonal KD-based baseline under the same heterogeneity
axes, useful for comparing dual-teacher (FLwF-2) and clustered teacher
(FedKDC-CL) mechanisms in future unified budgets.

%==================================================================
\appendix
\section{Full Results Table}
\label{app:results}
%==================================================================

Table~\ref{tab:full_results} reports the complete per-configuration
summary. ``Acc'' is final-round global test accuracy ($\%$),
``Comm.'' is cumulative communication cost in MB (full-weight, both
directions, FP32) and ``Forgetting'' is for Task~1. Convergence-round
is omitted for the reason given in Section~\ref{sec:analysis}.

% Full per-configuration results table (Appendix).
\begin{table*}[t]
\centering
\caption{Full experimental results summary across all 90 configurations.}
\label{tab:full_results}
\renewcommand{\arraystretch}{1.1}
\small
\begin{tabular}{llcccccc}
\toprule
\textbf{Method} & \textbf{Dataset} & \textbf{$K$}
 & \textbf{Partition} & \textbf{Rounds}
 & \textbf{Acc (\%)} & \textbf{Comm.\ (MB)} & \textbf{Forgetting} \\
\midrule
FedAvg & MNIST    &  10 & IID       & 20 &  48.3 &   47.8 & 1.00 \\
FedAvg & MNIST    &  10 & Dir(0.01) & 20 &  10.0 &   47.8 & 0.41 \\
FedAvg & MNIST    &  10 & Dir(0.1)  & 20 &  47.6 &   47.8 & 0.99 \\
FedAvg & MNIST    &  10 & Dir(0.5)  & 20 &  48.2 &   47.8 & 1.00 \\
FedAvg & MNIST    &  10 & Dir(1.0)  & 20 &  48.3 &   47.8 & 1.00 \\
FedAvg & MNIST    &  50 & IID       & 20 &  48.0 &  239.1 & 0.99 \\
FedAvg & MNIST    &  50 & Dir(0.01) & 20 &   9.7 &  239.1 & 0.24 \\
FedAvg & MNIST    &  50 & Dir(0.1)  & 20 &  44.8 &  239.1 & 0.97 \\
FedAvg & MNIST    &  50 & Dir(0.5)  & 20 &  47.9 &  239.1 & 0.99 \\
FedAvg & MNIST    &  50 & Dir(1.0)  & 20 &  47.9 &  239.1 & 0.99 \\
FedAvg & MNIST    & 100 & IID       & 20 &  47.8 &  478.3 & 0.99 \\
FedAvg & MNIST    & 100 & Dir(0.01) & 20 &  30.9 &  478.3 & 0.54 \\
FedAvg & MNIST    & 100 & Dir(0.1)  & 20 &  46.6 &  478.3 & 0.98 \\
FedAvg & MNIST    & 100 & Dir(0.5)  & 20 &  47.8 &  478.3 & 0.99 \\
FedAvg & MNIST    & 100 & Dir(1.0)  & 20 &  47.6 &  478.3 & 0.99 \\
\midrule
FedAvg & FMNIST   &  10 & IID       & 20 &  48.9 &   47.8 & 0.94 \\
FedAvg & FMNIST   &  10 & Dir(0.01) & 20 &  20.1 &   47.8 & 0.60 \\
FedAvg & FMNIST   &  10 & Dir(0.1)  & 20 &  47.6 &   47.8 & 0.83 \\
FedAvg & FMNIST   &  10 & Dir(0.5)  & 20 &  48.6 &   47.8 & 0.93 \\
FedAvg & FMNIST   &  10 & Dir(1.0)  & 20 &  48.7 &   47.8 & 0.93 \\
FedAvg & FMNIST   &  50 & IID       & 20 &  48.5 &  239.1 & 0.91 \\
FedAvg & FMNIST   &  50 & Dir(0.01) & 20 &  30.5 &  239.1 & 0.51 \\
FedAvg & FMNIST   &  50 & Dir(0.1)  & 20 &  47.9 &  239.1 & 0.81 \\
FedAvg & FMNIST   &  50 & Dir(0.5)  & 20 &  48.2 &  239.1 & 0.91 \\
FedAvg & FMNIST   &  50 & Dir(1.0)  & 20 &  48.3 &  239.1 & 0.91 \\
FedAvg & FMNIST   & 100 & IID       & 20 &  48.0 &  478.3 & 0.90 \\
FedAvg & FMNIST   & 100 & Dir(0.01) & 20 &  31.1 &  478.3 & 0.57 \\
FedAvg & FMNIST   & 100 & Dir(0.1)  & 20 &  46.7 &  478.3 & 0.82 \\
FedAvg & FMNIST   & 100 & Dir(0.5)  & 20 &  47.9 &  478.3 & 0.90 \\
FedAvg & FMNIST   & 100 & Dir(1.0)  & 20 &  48.0 &  478.3 & 0.90 \\
\midrule
FedAvg & CIFAR-10 &  10 & IID       & 20 &  43.2 &   48.3 & 0.80 \\
FedAvg & CIFAR-10 &  10 & Dir(0.01) & 20 &  11.5 &   48.3 & 0.45 \\
FedAvg & CIFAR-10 &  10 & Dir(0.1)  & 20 &  36.3 &   48.3 & 0.60 \\
FedAvg & CIFAR-10 &  10 & Dir(0.5)  & 20 &  40.7 &   48.3 & 0.79 \\
FedAvg & CIFAR-10 &  10 & Dir(1.0)  & 20 &  42.3 &   48.3 & 0.78 \\
FedAvg & CIFAR-10 &  50 & IID       & 20 &  38.8 &  241.4 & 0.74 \\
FedAvg & CIFAR-10 &  50 & Dir(0.01) & 20 &  11.8 &  241.4 & 0.32 \\
FedAvg & CIFAR-10 &  50 & Dir(0.1)  & 20 &  36.2 &  241.4 & 0.67 \\
FedAvg & CIFAR-10 &  50 & Dir(0.5)  & 20 &  35.0 &  241.4 & 0.73 \\
FedAvg & CIFAR-10 &  50 & Dir(1.0)  & 20 &  37.7 &  241.4 & 0.72 \\
FedAvg & CIFAR-10 & 100 & IID       & 20 &  37.4 &  482.9 & 0.69 \\
FedAvg & CIFAR-10 & 100 & Dir(0.01) & 20 &  11.9 &  482.9 & 0.30 \\
FedAvg & CIFAR-10 & 100 & Dir(0.1)  & 20 &  33.4 &  482.9 & 0.62 \\
FedAvg & CIFAR-10 & 100 & Dir(0.5)  & 20 &  36.7 &  482.9 & 0.67 \\
FedAvg & CIFAR-10 & 100 & Dir(1.0)  & 20 &  36.7 &  482.9 & 0.67 \\
\midrule
FLwF-2 & MNIST    &  10 & IID       & 20 &  51.3 &   47.8 & \textbf{0.00} \\
FLwF-2 & MNIST    &  10 & Dir(0.01) & 20 &   9.8 &   47.8 & 0.15 \\
FLwF-2 & MNIST    &  10 & Dir(0.1)  & 20 &  50.5 &   47.8 & \textbf{0.01} \\
FLwF-2 & MNIST    &  10 & Dir(0.5)  & 20 &  51.2 &   47.8 & \textbf{0.00} \\
FLwF-2 & MNIST    &  10 & Dir(1.0)  & 20 &  51.2 &   47.8 & \textbf{0.00} \\
FLwF-2 & MNIST    &  50 & IID       & 20 &  51.1 &  239.1 & \textbf{0.00} \\
FLwF-2 & MNIST    &  50 & Dir(0.01) & 20 &   9.8 &  239.1 & 0.20 \\
FLwF-2 & MNIST    &  50 & Dir(0.1)  & 20 &  50.6 &  239.1 & \textbf{0.00} \\
FLwF-2 & MNIST    &  50 & Dir(0.5)  & 20 &  51.0 &  239.1 & \textbf{0.00} \\
FLwF-2 & MNIST    &  50 & Dir(1.0)  & 20 &  51.1 &  239.1 & \textbf{0.00} \\
FLwF-2 & MNIST    & 100 & IID       & 20 &  50.9 &  478.3 & \textbf{0.00} \\
FLwF-2 & MNIST    & 100 & Dir(0.01) & 20 &  30.0 &  478.3 & \textbf{0.00} \\
FLwF-2 & MNIST    & 100 & Dir(0.1)  & 20 &  50.4 &  478.3 & \textbf{0.00} \\
FLwF-2 & MNIST    & 100 & Dir(0.5)  & 20 &  50.8 &  478.3 & \textbf{0.00} \\
FLwF-2 & MNIST    & 100 & Dir(1.0)  & 20 &  50.9 &  478.3 & \textbf{0.00} \\
\midrule
FLwF-2 & FMNIST   &  10 & IID       & 20 &  46.7 &   47.8 & \textbf{0.00} \\
FLwF-2 & FMNIST   &  10 & Dir(0.01) & 20 &  32.3 &   47.8 & \textbf{0.00} \\
FLwF-2 & FMNIST   &  10 & Dir(0.1)  & 20 &  37.0 &   47.8 & 0.09 \\
FLwF-2 & FMNIST   &  10 & Dir(0.5)  & 20 &  46.5 &   47.8 & \textbf{0.00} \\
FLwF-2 & FMNIST   &  10 & Dir(1.0)  & 20 &  46.1 &   47.8 & \textbf{0.01} \\
FLwF-2 & FMNIST   &  50 & IID       & 20 &  45.7 &  239.1 & \textbf{0.00} \\
FLwF-2 & FMNIST   &  50 & Dir(0.01) & 20 &  11.2 &  239.1 & 0.15 \\
FLwF-2 & FMNIST   &  50 & Dir(0.1)  & 20 &  40.8 &  239.1 & \textbf{0.00} \\
FLwF-2 & FMNIST   &  50 & Dir(0.5)  & 20 &  45.3 &  239.1 & \textbf{0.00} \\
FLwF-2 & FMNIST   &  50 & Dir(1.0)  & 20 &  45.5 &  239.1 & \textbf{0.00} \\
FLwF-2 & FMNIST   & 100 & IID       & 20 &  44.9 &  478.3 & \textbf{0.00} \\
FLwF-2 & FMNIST   & 100 & Dir(0.01) & 20 &  27.0 &  478.3 & \textbf{0.00} \\
FLwF-2 & FMNIST   & 100 & Dir(0.1)  & 20 &  34.5 &  478.3 & 0.13 \\
FLwF-2 & FMNIST   & 100 & Dir(0.5)  & 20 &  45.2 &  478.3 & \textbf{0.00} \\
FLwF-2 & FMNIST   & 100 & Dir(1.0)  & 20 &  44.9 &  478.3 & \textbf{0.00} \\
\midrule
FLwF-2 & CIFAR-10 &  10 & IID       & 20 &  40.1 &   48.3 & \textbf{0.00} \\
FLwF-2 & CIFAR-10 &  10 & Dir(0.01) & 20 &  10.0 &   48.3 & 0.26 \\
FLwF-2 & CIFAR-10 &  10 & Dir(0.1)  & 20 &  27.2 &   48.3 & 0.06 \\
FLwF-2 & CIFAR-10 &  10 & Dir(0.5)  & 20 &  39.3 &   48.3 & \textbf{0.00} \\
FLwF-2 & CIFAR-10 &  10 & Dir(1.0)  & 20 &  38.3 &   48.3 & \textbf{0.01} \\
FLwF-2 & CIFAR-10 &  50 & IID       & 20 &  36.7 &  241.4 & \textbf{0.00} \\
FLwF-2 & CIFAR-10 &  50 & Dir(0.01) & 20 &  11.1 &  241.4 & 0.07 \\
FLwF-2 & CIFAR-10 &  50 & Dir(0.1)  & 20 &  28.9 &  241.4 & 0.09 \\
FLwF-2 & CIFAR-10 &  50 & Dir(0.5)  & 20 &  34.4 &  241.4 & \textbf{0.04} \\
FLwF-2 & CIFAR-10 &  50 & Dir(1.0)  & 20 &  36.0 &  241.4 & \textbf{0.00} \\
FLwF-2 & CIFAR-10 & 100 & IID       & 20 &  34.1 &  482.9 & \textbf{0.00} \\
FLwF-2 & CIFAR-10 & 100 & Dir(0.01) & 20 &   9.9 &  482.9 & 0.10 \\
FLwF-2 & CIFAR-10 & 100 & Dir(0.1)  & 20 &  31.8 &  482.9 & \textbf{0.00} \\
FLwF-2 & CIFAR-10 & 100 & Dir(0.5)  & 20 &  33.0 &  482.9 & \textbf{0.00} \\
FLwF-2 & CIFAR-10 & 100 & Dir(1.0)  & 20 &  33.6 &  482.9 & \textbf{0.00} \\
\bottomrule
\end{tabular}
\end{table*}


%==================================================================
\bibliographystyle{IEEEtran}
\begin{thebibliography}{99}

\bibitem{fedavg}
B. McMahan, E. Moore, D. Ramage, S. Hampson, and B. A. y Arcas,
``Communication-efficient learning of deep networks from decentralized data,''
\textit{Proc. Artif. Intell. Statist.}, pp.\ 1273--1282, 2017.

\bibitem{flwf}
A. Usmanova, F. Portet, P. Lalanda, and G. Vega,
``Federated continual learning through distillation in pervasive computing,''
\textit{arXiv:2207.08181}, 2022.

\bibitem{hinton_kd}
G. Hinton, O. Vinyals, and J. Dean,
``Distilling the knowledge in a neural network,''
\textit{NIPS Deep Learning Workshop}, 2015.

\bibitem{forgetting}
J. Kirkpatrick et al.,
``Overcoming catastrophic forgetting in neural networks,''
\textit{Proc. Natl. Acad. Sci.}, vol.\ 114, no.\ 13, pp.\ 3521--3526, 2017.

\bibitem{flower}
D. J. Beutel et al.,
``Flower: A friendly federated learning research framework,''
\textit{arXiv:2007.14390}, 2020.

\bibitem{lwf}
Z. Li and D. Hoiem,
``Learning without forgetting,''
\textit{IEEE Trans. Pattern Anal. Mach. Intell.}, vol.\ 40, pp.\ 2935--2947, 2017.

\bibitem{feddf}
T. Lin, L. Kong, S. U. Stich, and M. Jaggi,
``Ensemble distillation for robust model fusion in federated learning,''
\textit{Proc. NeurIPS}, vol.\ 33, pp.\ 2351--2363, 2020.

\bibitem{fedmd}
D. Li and J. Wang,
``FedMD: Heterogenous federated learning via model distillation,''
\textit{NeurIPS Workshop}, 2019.

\bibitem{fedprox}
T. Li, A. K. Sahu, M. Zaheer, M. Sanjabi, A. Talwalkar, and V. Smith,
``Federated optimization in heterogeneous networks,''
\textit{Proc. MLSys}, vol.\ 2, pp.\ 429--450, 2020.

\bibitem{dirichlet}
K.-W. Ng, G.-L. Tian, and M.-L. Tang,
``Dirichlet and related distributions: Theory, methods and applications,''
Wiley Series in Probability and Statistics, Wiley, 2011.

\end{thebibliography}

\end{document}
